\chapter{运行与维护}

本章阐述系统的运行环境配置、性能管理策略、可靠性保障机制以及日常运维检查要点。

\section{运行环境}

系统的运行时组件栈自上而下由以下部分构成：

\begin{description}
  \item[浏览器端] 用户通过现代浏览器访问单页应用，前端以静态资源形式部署，降低运行环境复杂度。
  \item[网关层] 生产环境由 Nginx 统一提供静态文件服务和反向代理，并负责 HTTPS、压缩和缓存等通用能力。
  \item[后端服务] 后端应用以内网服务形式运行，统一提供认证、导入、分类、分析、组织和报表等业务能力，并由进程管理工具保障持续运行。
  \item[数据与任务组件] 生产环境使用 PostgreSQL 保存业务数据，Redis 支撑异步任务与缓存；轻量部署可采用 SQLite 和同步任务模式降低部署成本。
  \item[模型服务] 分类模型既可接入外部大模型服务，也可切换至本地模型或规则兜底。模型请求由重试机制统一调度，避免外部服务波动直接影响用户操作。
\end{description}

文件存储方面，用户上传账单和生成报表分目录保存，并按用户维度隔离。运维上应定期清理过期上传文件，控制磁盘占用，同时保留可追溯的报表输出记录。

\section{性能管理}

下表汇总了各核心组件的已知性能瓶颈、当前缓解措施及未来优化方向：

\begin{table}[H]
  \centering
  \small
  \caption{性能管理汇总}
  \label{tab:performance}
  \begin{tabular}{@{}p{2.2cm} p{2.8cm} p{3.8cm} p{4.0cm}@{}}
    \toprule
    \textbf{组件} & \textbf{瓶颈描述} & \textbf{当前缓解措施} & \textbf{未来优化方向} \\
    \midrule
    账单导入解析 & 读取与转换大文件时占用 I/O 和内存 & 按文件分批处理，限制扫描范围，避免长事务 & 大文件异步处理，并行解析与批量写入 \\
    \hline
    模型语义分类 & 外部模型延迟和可用性波动明显 & 统一重试池管理外部调用；缓存重复分类结果 & 引入语义相似缓存；支持批量分类减少网络往返 \\
    \hline
    Dashboard 聚合 & 趋势与分类统计需扫描交易记录 & 使用索引和日期范围减少扫描量 & 引入物化视图或预聚合表，定时刷新统计 \\
    \hline
    消费人格计算 & Shannon 指数与余弦相似度涉及全量交易扫描 & 当前规模下（单用户数千笔）耗时百毫秒级，可接受 & 引入增量计算，新交易更新统计量而非全量扫描 \\
    \hline
    PDF 报表生成 & 多页报表渲染消耗 CPU 时间 & 典型报表秒级生成；异步任务避免阻塞用户请求 & 超大型报表分页生成后合并 \\
    \hline
    前端加载性能 & 组件库和图表库体积较大 & 生产构建启用按需优化与压缩 & 路由级代码分割，按页面懒加载图表能力 \\
    \bottomrule
  \end{tabular}
\end{table}

\section{可靠性保障}

系统通过以下五项机制保障数据处理的可靠性：

\begin{enumerate}
  \item \textbf{重复交易控制}。每笔交易入库前都会生成稳定指纹，并按用户维度进行唯一性校验。同一文件重复导入时，系统自动跳过重复记录，避免统计结果膨胀。
  \item \textbf{统一重试池}。外部模型调用与用户请求解耦，失败请求进入后台重试流程。重试达到上限后，交易被标记为需要管理员关注，避免失败任务持续占用处理资源。
  \item \textbf{导入批次状态机}。导入任务按等待、处理中、完成、部分失败和失败等状态流转，页面通过进度条和状态标签向用户解释当前处理结果。
  \item \textbf{报表任务状态机}。报表生成作为异步任务执行，用户可在生成完成后下载，避免大报表生成阻塞页面操作。
  \item \textbf{权限校验}。系统在交易查询、分类修正、批次删除、报表下载和组织管理等操作前校验身份与资源归属，防止跨用户或跨组织访问。
\end{enumerate}

\section{运维检查}

日常运维中建议按照以下检查清单定期审查系统状态：

\begin{table}[H]
  \centering
  \small
  \caption{运维检查清单}
  \label{tab:ops-checklist}
  \begin{tabular}{@{}p{2.4cm} p{5.0cm} p{5.4cm}@{}}
    \toprule
    \textbf{检查项} & \textbf{检查方式} & \textbf{预期结果与说明} \\
    \midrule
    日志监控 & 定期查看访问日志与错误日志 & 关注 5xx 错误频率、超时日志和异常堆栈，定位异常用户流程 \\
    \hline
    健康检查 & 使用健康检查端点或监控系统探活 & 服务应持续返回正常状态，异常时及时告警 \\
    \hline
    数据库备份 & 配置自动备份并定期演练恢复 & 生产环境建议每日备份，保留最近 7 天并验证可恢复性 \\
    \hline
    密钥轮换 & 按计划更新认证密钥和模型访问凭证 & 轮换后旧凭证失效，建议低峰期执行并保留回滚方案 \\
    \hline
    依赖安全审计 & 定期检查前后端依赖漏洞报告 & 关注高危漏洞，及时更新受影响的依赖版本 \\
    \hline
    上传文件控制 & 检查文件大小、类型限制和存储目录 & 确保符合安全策略，定期清理过期上传文件 \\
    \hline
    权限回归测试 & 修改认证或组织逻辑后运行相关测试 & 确保用户隔离和家庭角色权限未退化 \\
    \hline
    重试队列监控 & 管理员查看重试状态与失败数量 & 若失败数量持续增长，应检查外部模型服务可用性 \\
    \hline
    磁盘空间检查 & 定期检查上传文件和报表目录容量 & 不应超过预期；建议定时清理过期上传文件 \\
    \bottomrule
  \end{tabular}
\end{table}

此外，建议在生产服务器上配置以下监控：

\begin{itemize}
  \item \textbf{进程存活监控}：确保前端网关、后端服务和后台任务异常退出后能够自动重启。
  \item \textbf{内存与 CPU 监控}：关注导入、报表生成和模型调用期间的资源尖峰，及时发现内存泄漏或计算瓶颈。
  \item \textbf{磁盘 I/O 监控}：关注上传文件、报表输出和数据库写入对磁盘的影响，必要时切换至更适合并发写入的数据库方案。
  \item \textbf{外部模型可用性告警}：若重试失败数量持续超过阈值，应触发告警通知运维人员检查模型服务可用性与配额状态。
\end{itemize}
